%! Author = denis
%! Date = 11/08/2026

\begin{frame}{Identidades lógicas: Lei da Idempotente}
    \begin{block}{Definição}
        Repetir uma mesma variável em uma operação lógica não altera
        o seu resultado.
    \end{block}

    \vspace{0.3cm}

    \begin{columns}[T]

        % Coluna OR
        \begin{column}{0.5\textwidth}
            \begin{block}{Idempotência OR}
                \[
                    \boxed{A + A = A}
                \]
                \[
                    0 + 0 = 0
                \]
                \[
                    1 + 1 = 1
                \]

                Portanto, repetir $A$ não altera a saída.
            \end{block}
        \end{column}

        % Coluna AND
        \begin{column}{0.5\textwidth}
            \begin{block}{Idempotência AND}
                \[
                    \boxed{A \cdot A = A}
                \]
                \[
                    0 \cdot 0 = 0
                \]
                \[
                    1 \cdot 1 = 1
                \]

                Portanto, repetir $A$ não altera a saída.
            \end{block}
        \end{column}

    \end{columns}

    \begin{alertblock}{Interpretação}
        Na Álgebra Booleana, repetir a mesma variável em uma operação
        \textbf{OR} ou \textbf{AND} é redundante e pode ser eliminado.
    \end{alertblock}
\end{frame}